\documentclass[14pt,a4paper]{extarticle}

\usepackage[T2A]{fontenc}
\usepackage[utf8]{inputenc}
\usepackage[russian]{babel}
\usepackage{cmap}
\usepackage{geometry}
\usepackage{array}
\usepackage{longtable}
\usepackage{booktabs}
\usepackage{xcolor}
\usepackage{ragged2e}
\usepackage{seqsplit}
\usepackage{listings}
\usepackage{enumitem}
\usepackage{hyperref}
\usepackage{tikz}
\usetikzlibrary{arrows.meta,positioning,shapes.geometric}

\geometry{
    left=30mm,
    right=15mm,
    top=20mm,
    bottom=20mm
}

\hypersetup{
    colorlinks=true,
    linkcolor=black,
    urlcolor=blue
}

\setlength{\parindent}{1.25cm}
\setlength{\parskip}{0.2em}
\renewcommand{\arraystretch}{1.35}
\setlength{\tabcolsep}{4pt}
\sloppy

\newcolumntype{L}[1]{>{\RaggedRight\arraybackslash}p{#1}}

\lstdefinestyle{codestyle}{
    basicstyle=\ttfamily\small,
    breaklines=true,
    frame=single,
    columns=fullflexible,
    keepspaces=true,
    showstringspaces=false,
    backgroundcolor=\color{gray!5}
}

\newcommand{\code}[1]{\texttt{\detokenize{#1}}}
\newcommand{\pathcode}[1]{{\footnotesize\nolinkurl{#1}}}

\begin{document}

\begin{titlepage}
    \begin{center}
        САНКТ-ПЕТЕРБУРГСКИЙ НАЦИОНАЛЬНЫЙ\\
        ИССЛЕДОВАТЕЛЬСКИЙ УНИВЕРСИТЕТ ИТМО

        \vspace{1.5cm}

        Дисциплина: Бэк-энд разработка

        \vspace{1.5cm}

        \textbf{Отчет}

        \vspace{0.5cm}

        Лабораторная работа 3

        \vspace{0.5cm}

        \textbf{Контейнеризация написанного приложения средствами Docker}

        \vfill

        \begin{flushright}
            Выполнил:\\
            Цатинян Артём\\
            Группа БР1.2

            \vspace{0.8cm}

            Проверил:\\
            Добряков Д. И.
        \end{flushright}

        \vfill

        Санкт-Петербург\\
        2026 г.
    \end{center}
\end{titlepage}

\section{Задача}

Тема проекта: приложение для бронирования столиков в ресторанах.

В рамках лабораторной работы требовалось контейнеризовать микросервисное backend-приложение: реализовать \code{Dockerfile} для каждого сервиса, написать общий \code{docker-compose.yml} и настроить сетевое взаимодействие между сервисами.

\section{Ход работы}

Контейнеризация выполнялась для микросервисной версии приложения, реализованной в рамках ЛР2. Проект представляет собой Gradle multi-module приложение на Kotlin и Spring Boot. В состав системы входят четыре бизнес-сервиса и две инфраструктурные зависимости: PostgreSQL и RabbitMQ.

\subsection{Состав контейнеризуемого приложения}

\begin{longtable}{|L{0.28\textwidth}|L{0.42\textwidth}|L{0.16\textwidth}|}
\hline
\textbf{Сервис} & \textbf{Назначение} & \textbf{Порт} \\
\hline
\code{identity-service} & Регистрация, авторизация, выдача JWT, получение данных текущего пользователя & 8081 \\
\hline
\code{catalog-service} & Каталог ресторанов, кухни, меню, столики, агрегированные рейтинги & 8082 \\
\hline
\code{booking-service} & Создание и просмотр бронирований, проверка доступности столика & 8083 \\
\hline
\code{review-service} & Создание и просмотр отзывов, публикация событий об отзывах & 8084 \\
\hline
\code{postgres} & Общая PostgreSQL база для локального запуска, схемы разделены по сервисам & 5432 \\
\hline
\code{rabbitmq} & Брокер сообщений для асинхронного взаимодействия & 5672, 15672 \\
\hline
\end{longtable}

\subsection{Dockerfile для сервисов}

Для каждого бизнес-сервиса был создан отдельный \code{Dockerfile}. Все файлы построены по одной схеме:

\begin{enumerate}[nosep]
    \item На первом этапе используется образ \code{eclipse-temurin:21-jdk}.
    \item В контейнер копируется весь multi-module проект.
    \item Gradle собирает bootJar только нужного сервиса.
    \item На втором этапе используется более легкий runtime-образ \code{eclipse-temurin:21-jre}.
    \item В итоговый образ копируется только готовый jar-файл сервиса.
\end{enumerate}

Пример общей логики Dockerfile:

\begin{lstlisting}[style=codestyle]
FROM eclipse-temurin:21-jdk AS build
WORKDIR /workspace
COPY . .
RUN chmod +x ./gradlew \
    && ./gradlew --console=plain :booking-service:bootJar \
       -x liquibaseUpdate -x generateJooq --no-daemon

FROM eclipse-temurin:21-jre
WORKDIR /app
COPY --from=build /workspace/booking-service/build/libs/booking-service.jar app.jar
EXPOSE 8083
ENTRYPOINT ["java", "-jar", "/app/app.jar"]
\end{lstlisting}

Генерация jOOQ и Liquibase update не запускаются внутри Docker build. Сгенерированные jOOQ-классы уже находятся в проекте, а миграции Liquibase применяются при старте каждого Spring Boot сервиса.

\subsection{Оптимизация контекста сборки}

Для уменьшения Docker build context добавлен файл \code{.dockerignore}. Из контекста исключены локальные Gradle-кэши, build-директории, временные файлы и PDF/TeX артефакты отчетов. Это ускоряет сборку образов и не тащит внутрь контейнеров лишние локальные файлы.

\subsection{Общий docker-compose.yml}

Для запуска всей системы добавлен общий \code{docker-compose.yml}. Он поднимает PostgreSQL, RabbitMQ и все четыре микросервиса в одной Docker-сети. Благодаря этому сервисы обращаются друг к другу не через \code{localhost}, а по DNS-именам контейнеров.

\begin{center}
\begin{tikzpicture}[
    node distance=1.9cm,
    service/.style={rectangle, draw, rounded corners, align=center, minimum width=3.2cm, minimum height=0.9cm},
    infra/.style={cylinder, draw, shape border rotate=90, aspect=0.25, align=center, minimum width=2.8cm, minimum height=1cm},
    arrow/.style={-{Latex[length=3mm]}, thick}
]
\node[service] (identity) {identity-service\\8081};
\node[service, right=2.4cm of identity] (catalog) {catalog-service\\8082};
\node[service, below=1.6cm of identity] (booking) {booking-service\\8083};
\node[service, right=2.4cm of booking] (review) {review-service\\8084};
\node[infra, below=1.5cm of booking] (postgres) {PostgreSQL\\5432};
\node[infra, below=1.5cm of review] (rabbit) {RabbitMQ\\5672/15672};

\draw[arrow] (booking) -- node[left]{auth} (identity);
\draw[arrow] (booking) -- node[above]{catalog context} (catalog);
\draw[arrow] (review) -- node[right]{auth} (identity);
\draw[arrow] (review) -- node[above]{booking context} (booking);
\draw[arrow] (review) -- node[right]{events} (rabbit);
\draw[arrow] (rabbit) -- node[right]{events} (catalog);
\draw[arrow] (identity) -- (postgres);
\draw[arrow] (catalog) -- (postgres);
\draw[arrow] (booking) -- (postgres);
\draw[arrow] (review) -- (postgres);
\end{tikzpicture}
\end{center}

\subsection{Сетевое взаимодействие}

В Docker Compose все сервисы находятся в одной сети. Поэтому для межсервисных вызовов используются внутренние адреса:

\begin{longtable}{|L{0.35\textwidth}|L{0.45\textwidth}|}
\hline
\textbf{Назначение} & \textbf{Адрес внутри Docker-сети} \\
\hline
PostgreSQL & \pathcode{jdbc:postgresql://postgres:5432/storage_local} \\
\hline
RabbitMQ & \code{rabbitmq:5672} \\
\hline
Identity Service & \code{http://identity-service:8081} \\
\hline
Catalog Service & \code{http://catalog-service:8082} \\
\hline
Booking Service & \code{http://booking-service:8083} \\
\hline
\end{longtable}

Например, \code{booking-service} получает адреса зависимых сервисов через переменные окружения:

\begin{lstlisting}[style=codestyle]
SERVICES_CATALOG_BASE_URL: http://catalog-service:8082
SERVICES_IDENTITY_BASE_URL: http://identity-service:8081
\end{lstlisting}

\code{review-service} аналогично получает адреса \code{booking-service} и \code{identity-service}, а также параметры подключения к RabbitMQ.

\subsection{Порядок старта}

Для инфраструктуры настроены healthcheck-и. PostgreSQL считается готовым после успешного \code{pg_isready}, RabbitMQ --- после успешного \code{rabbitmq-diagnostics ping}. Бизнес-сервисы запускаются после готовности нужной инфраструктуры.

\begin{longtable}{|L{0.28\textwidth}|L{0.52\textwidth}|}
\hline
\textbf{Сервис} & \textbf{Зависимости при запуске} \\
\hline
\code{identity-service} & \code{postgres} \\
\hline
\code{catalog-service} & \code{postgres}, \code{rabbitmq} \\
\hline
\code{booking-service} & \code{postgres}, \code{identity-service}, \code{catalog-service} \\
\hline
\code{review-service} & \code{postgres}, \code{rabbitmq}, \code{identity-service}, \code{booking-service} \\
\hline
\end{longtable}

\subsection{Запуск и проверка}

Для запуска всей системы используется команда:

\begin{lstlisting}[style=codestyle]
docker-compose up --build -d
\end{lstlisting}

Если установлен новый Docker Compose plugin, можно использовать вариант:

\begin{lstlisting}[style=codestyle]
docker compose up --build -d
\end{lstlisting}

После запуска доступны Swagger UI каждого сервиса:

\begin{longtable}{|L{0.28\textwidth}|L{0.55\textwidth}|}
\hline
\textbf{Сервис} & \textbf{Swagger UI} \\
\hline
\code{identity-service} & \pathcode{http://localhost:8081/swagger-ui/index.html} \\
\hline
\code{catalog-service} & \pathcode{http://localhost:8082/swagger-ui/index.html} \\
\hline
\code{booking-service} & \pathcode{http://localhost:8083/swagger-ui/index.html} \\
\hline
\code{review-service} & \pathcode{http://localhost:8084/swagger-ui/index.html} \\
\hline
\end{longtable}

Подключение к базе данных:

\begin{lstlisting}[style=codestyle]
host: localhost
port: 5432
database: storage_local
username: storage
password: storage
\end{lstlisting}

RabbitMQ Management UI:

\begin{lstlisting}[style=codestyle]
url: http://localhost:15672
username: storage
password: storage
\end{lstlisting}

Для остановки контейнеров используется команда:

\begin{lstlisting}[style=codestyle]
docker-compose down
\end{lstlisting}

Для полной очистки данных PostgreSQL:

\begin{lstlisting}[style=codestyle]
docker-compose down -v
\end{lstlisting}

\section{Вывод}

В рамках лабораторной работы была выполнена контейнеризация микросервисного backend-приложения. Для каждого сервиса создан отдельный Dockerfile, а общий \code{docker-compose.yml} поднимает всю систему вместе с PostgreSQL и RabbitMQ. Межсервисное взаимодействие настроено через внутренние DNS-имена Docker Compose, поэтому сервисы могут обращаться друг к другу внутри общей сети без ручной настройки адресов.

В результате приложение можно запускать одной командой, что упрощает демонстрацию на защите и приближает локальное окружение к реальному развертыванию микросервисной системы.

\end{document}
